\documentclass[11pt]{article}
\usepackage[a4paper,margin=1in]{geometry}
\usepackage{amsmath,amssymb,amsthm,microtype}
\newtheorem{theorem}{Theorem}[section]
\theoremstyle{remark}\newtheorem{remark}[theorem]{Remark}
\newcommand{\Tr}{\operatorname{Tr}}
\title{A Weighted Prefix--Tail Gluing Criterion for Regularized Determinants}
\author{Bin Wang}\date{July 2026}
\begin{document}\maketitle
\begin{abstract}
The new spectral-tail certificate and the deterministic numerator envelope
live on opposite sides of a missing prefix interface.  We formulate the
exact gluing theorem.  Split the logarithmic coefficient difference at a
moving order $m_\sigma$.  If the weighted prefix mismatch, the noisy
high-order tail, and the target high-order tail all vanish on a closed disk,
then the normalized determinant quotient converges uniformly there, with an
explicit exponential error bound.  RH-285 supplies the noisy spectral-tail
leaf when the noisy coefficient sequence is the modulus-complement trace,
and the deterministic envelope papers supply the target leaf.  The remaining
prefix has a typed two-error decomposition: total noisy trace versus
counterloop plus anchor, and noisy head versus counterloop.  RH-287 controls
only the first error unweighted and without a rate.  The criterion prevents
finite-order agreement from being promoted to a determinant theorem.
\end{abstract}

\section{Three budgets}
Let
\[
 F_\sigma(z)=\exp\left(-\sum_{n\ge2}\frac{b_{\sigma,n}}n z^n\right),
 \qquad
 F(z)=\exp\left(-\sum_{n\ge2}\frac{a_n}n z^n\right)
\]
be normalized logarithmic germs whose displayed logarithmic series converge
absolutely and uniformly on the closed disk $|z|\le R$.
For an integer $m_\sigma\ge3$, put
\begin{align*}
 P_\sigma(R)&=\sum_{2\le n<m_\sigma}
 \frac{|b_{\sigma,n}-a_n|R^n}{n},\\
 S_\sigma(R)&=\sum_{n\ge m_\sigma}
 \frac{|b_{\sigma,n}|R^n}{n},\\
 T_\sigma(R)&=\sum_{n\ge m_\sigma}
 \frac{|a_n|R^n}{n}.
\end{align*}

\section{Gluing theorem}
\begin{theorem}\label{thm:gluing}
If $P_\sigma(R)+S_\sigma(R)+T_\sigma(R)\to0$, then
\[
 \sup_{|z|\le R}\left|\frac{F_\sigma(z)}{F(z)}-1\right|
 \le
 \exp\{P_\sigma(R)+S_\sigma(R)+T_\sigma(R)\}-1
 \longrightarrow0.
\]
\end{theorem}
\begin{proof}
On $|z|\le R$, split the normalized logarithm
\[
 \log\frac{F_\sigma(z)}{F(z)}
 =-\sum_{2\le n<m_\sigma}\frac{(b_{\sigma,n}-a_n)z^n}{n}
  -\sum_{n\ge m_\sigma}\frac{b_{\sigma,n}z^n}{n}
  +\sum_{n\ge m_\sigma}\frac{a_nz^n}{n}.
\]
Its modulus is at most $P_\sigma+S_\sigma+T_\sigma$.  The inequality
$|e^w-1|\le e^{|w|}-1$ gives the result.
\end{proof}

\section{Typed spectral application}
Let $c_{\sigma,n}$ be the total noisy bulk trace, let
\[
 h_{\sigma,n}=\sum_{\mu\in H_\sigma}\mu^n,
 \qquad
 \tau_{\sigma,n}=c_{\sigma,n}-h_{\sigma,n}
                  =\Tr C_\sigma^n,
\]
and let $s_{k_\sigma,n}$ be the finite-radius counterloop moment.  The target
anchor is $a_n$.  Define the two typed errors
\[
 e_{\sigma,n}=c_{\sigma,n}-s_{k_\sigma,n}-a_n,
 \qquad
 d_{\sigma,n}=h_{\sigma,n}-s_{k_\sigma,n}.
\]
Then the coefficient needed for the spectral complement satisfies the exact
identity
\begin{equation}\label{eq:typed}
 \tau_{\sigma,n}-a_n=e_{\sigma,n}-d_{\sigma,n}.
\end{equation}
Consequently its prefix budget obeys
\[
 P_\sigma(R)\le E_\sigma(R)+D_\sigma(R),
\]
where
\[
 E_\sigma(R)=\sum_{2\le n<m_\sigma}
 \frac{|e_{\sigma,n}|R^n}{n},\qquad
 D_\sigma(R)=\sum_{2\le n<m_\sigma}
 \frac{|d_{\sigma,n}|R^n}{n}.
\]
Thus one may prove the missing prefix either directly for
$\tau_{\sigma,n}-a_n$, or by
proving both typed weighted estimates.  Head-to-counterloop transport alone
does not imply the complement-to-anchor estimate.

\section{Current leaf audit}
In the gluing theorem set $b_{\sigma,n}=\tau_{\sigma,n}$.  The target trace
envelope and coefficient-radius law give $T_\sigma(R)\to0$ on the certified
strict target disk.  RH-282 and RH-285 give $S_\sigma(R)\to0$ for this
projection-free modulus complement after the logarithmic block clock.  The
unresolved term is the direct complement-to-anchor prefix $P_\sigma(R)$.

RH-287 proves an unweighted maximum for $e_{\sigma,n}$ on some growing
prefix, but its only automatic weighted consequence is
\[
 E_\sigma^{(h)}(R):=
 \sum_{2\le n\le h_\sigma}\frac{|e_{\sigma,n}|R^n}{n}
 \le\varepsilon_\sigma
 \sum_{2\le n\le h_\sigma}\frac{R^n}{n}.
\]
The clock $h_\sigma$ need not reach the spectral tail clock $m_\sigma$, and
for $R>1$ even the displayed partial bound need not vanish.  No current
theorem supplies the weighted
$D_\sigma(R)$ term either.  A direct $P_\sigma\to0$ theorem, or both typed
weighted estimates in \eqref{eq:typed}, is still required.

\begin{remark}
The theorem is an assembly criterion, not a claim that its three leaves are
all present.  Gate A and Gates B--E remain false/open.  It does not construct
a Hilbert--Polya operator or identify any arithmetic divisor.
\end{remark}
\end{document}
